\section{K-means 在交替下降中制造簇}
\label{sec:13-Machine-Learning/05-Clustering-and-Data-Geometry/01}

手里一堆没有标签的点 \(x_1, \ldots, x_n \in \R^p\)，你想把它们分成有意义的组。麻烦在于``有意义''本身没有数学定义——十个人可能划出十种分法，各说各的理。K-means 的做法不是试图定义普适的``有意义''，而是先选定一项明确的、可计算的好坏标准：好的分组应当让每个点尽可能接近它所属簇的 Euclidean 中心。

这不是聚类问题的唯一答案，但它是一个清楚到可以被计算、被批评、被改进的起点。

\subsection{两个未知量互相定义彼此}

问题里藏着两个彼此依赖的未知量：每个点属于哪个簇（分配 \(z_i\)），以及每个簇的中心在哪里（\(\mu_k\)）。单独问任何一个似乎都可以回答——只要另一个已经知道。

如果中心 \(\mu_1, \ldots, \mu_K\) 已知，分配问题退化到每个点独立决策：去最近的中心。

\[
z_i \leftarrow \argmin_{k \in \{1,\ldots,K\}} \norm{x_i - \mu_k}^2.
\]

如果分配已知，每个中心的最佳位置是它手下所有点的均值——这是平方误差下唯一最小化簇内距离的位置。

\[
\mu_k = \frac{1}{n_k} \sum_{i: z_i = k} x_i,
\qquad n_k = |\{i: z_i = k\}|.
\]

两者互相定义。于是自然出现一个迭代方案：随便猜一组中心，分配；用分配结果重新算中心；再分配；反复。这个方案的完整写法是把两个步骤放进一个统一的目标函数：

\[
J(z, \mu) = \sum_{i=1}^{n} \norm{x_i - \mu_{z_i}}^2.
\]

固定 \(\mu\) 更新 \(z\)，\(J\) 不增（每个点换到最近中心，平方距离只可能减少或持平）。固定 \(z\) 更新 \(\mu\)，\(J\) 也不增（均值是平方误差的全局最小化解）。两个步骤交替，\(J\) 单调下降。

\subsection{尝试一次：从随机初始中心出发}

取 \(K=3\)，从数据里随机挑 3 个点作为初始中心 \(\mu_1^{(0)}, \mu_2^{(0)}, \mu_3^{(0)}\)。第一轮分配：每个点去找离自己最近的中心。簇的边界自然出现——两个中心连线的垂直平分面，在二维就是垂直平分线。第一轮更新：重新计算每个簇的均值，中心移动。第二轮分配：一些点换了阵营，边界位移。几轮之后，中心不再明显移动，\(J\) 趋于平稳。

这个流程叫 Lloyd 算法。分配步骤的搜索空间是离散的——每个点有 \(K\) 种选择——但有限。每次分配只可能改变分区的排列，而 \(n\) 个点到 \(K\) 个簇的分配方式总数是有限的（虽然巨大）。结合 \(J\) 单调下降，算法必在有限步内停止在某个不能再通过单步操作改进的状态。每个这样的状态是一个局部稳定分区，但不保证是全局最优分区。

\subsection{不同起点，不同终点}

现在做一件事：换一组随机初始中心，再跑一次。大概率你会得到不同的最终分区和不同的最终 \(J\) 值。这个现象不是 bug，而是目标函数本身多峰的必然结果。

考虑一个简单的失败场景：真实数据大致分为三个分离良好的群，但你随机初始化的两个中心不幸都落在了同一个群里。那个群被两个中心瓜分，远处的第三个群被剩下的一个中心草草覆盖。随后的均值更新中，瓜分同一群的两个中心各自被局部数据拖住，无法跳出去寻找远处的数据。Lloyd 算法不包含``跳出''的机制——每一步只允许向局部有利的方向移动。

K-means++ 初始化尝试降低这种糟糕事件的概率。它不是等概率随机挑选 \(K\) 个初始中心，而是让新中心更倾向于选择离已有中心较远的点。第一个中心随机选；第二个中心选离第一个远的点，概率正比于平方距离；第三个中心选离前两个都远的点；以此类推。数学可以证明，这种初始化使得最终 \(J\) 的期望不超过最优值的 \(O(\log K)\) 倍。

但概率降低不等于消除。可靠的实践是：使用多次独立初始化，比较各次的最终 \(J\) 值，并检查各次的分区是否高度重合。仅凭``选 \(J\) 最小的一次''仍不足以保证结构可靠——高维噪声或异常点也可能产生很低的目标值，但分出来的簇对应不到任何真实模式。

\subsection{平方距离写了一部隐藏的几何法规}

K-means 的分区边界来自``到最近中心''的竞争，在欧氏空间中产生的是 Voronoi 划分——每两个中心之间有一个垂直平分超平面，整体是一组凸多面体。这份几何偏好几样东西：

它偏好尺度相近的簇。一个大而松散的簇可能被切成两半，因为它的远端点离自身中心比离邻近小簇的中心更远。

它偏好球形的簇。两个月牙形区域即使视觉上各自连通，月牙的尖端离对方中心可能比离自己中心更近，结果月牙被拦腰切开。

它对密度差异敏感。一个高密度小簇的中心可能被周围低密度大簇的少量点拖走，所以标准化不是可选的预处理——它是你对各特征相对重要性的宣告。将一个坐标乘以 10，相当于宣布该坐标上的距离重要了 100 倍（因为平方）。不同的标准化会导致完全不同的分区。

类别变量不能直接塞进 K-means。把``城市 = \{北京, 上海, 广州\}''编码为 \(\{1, 2, 3\}\) 意味着``北京距上海的距离 = 1，距广州的距离 = 2''——这种距离关系是编码方案强加的，不是数据本身携带的。类别差异没有自然欧氏几何。

\subsection{空簇与异常点：目标函数不说话的时候}

某轮分配后，可能出现某个中心没有收到任何点。均值的分母 \(n_k = 0\)，更新无定义。这不是浮点运算异常——它是离散优化完全合法的状态。处理方式通常是在实现层面约定：重新随机初始化空簇的中心、从最大的簇中拆出一半给它，或在 \(K\) 的选择上退一步。数学定义没有义务为所有参数值提供有意义的输出。

另一个更常见的问题是异常点利用平方距离放大自身影响。一个远离所有簇的点，它的平方距离可能是普通点的一百倍。中心会被这个点拖向它，而大量正常点因此略微远离了中心——总平方误差可能还下降了，但簇的位置被一个点劫持了。

K-medoids 是对这个弱点的直接回应：限制中心必须是实际数据点（而非均值），并使用一般的距离度量替代平方欧氏距离。代价是每轮更新需要搜索所有点作为候选中心，计算成本从 \(O(n)\) 上升到 \(O(n^2)\)。选择 K-medoids 是在稳健性和计算效率之间的一次有意识的交换。

\subsection{选择 \(K\)：训练目标不会帮你}

\(K\) 增大，训练目标 \(J\) 必然不降——每个点可以选择更近的中心。到极端，\(K=n\) 时每个点自成中心，\(J=0\)。所以``选使 \(J\) 最小的 \(K\)''总是退化为选 \(K=n\)。不能靠训练目标来选簇数——这是所有聚类方法的共同困境，不唯独是 K-means 的问题。

Elbow 方法看 \(J\) 随 \(K\) 的下降曲线，寻``变缓的拐点''。Silhouette 分数衡量每个点是否更接近自身簇内的点而非邻近簇的点。稳定性方法检查分区是否在不同随机子样本下保持一致。概率模型准则（如 BIC）假设数据来自某个混合分布并惩罚参数个数。下游任务则直接用最终用途（如客户分群后的营销效果）评判簇的质量。

不同的准则量度不同的性质，可能给出不同的答案。它们之间的分歧不应被视为``找到正确答案''的失败，而应被理解为：``簇数''本身在数据中可能就不唯一。一个数据集可以同时被合理地分为宏观的 3 组和微观的 7 组——两种分法揭示的是同一数据中不同尺度的结构。

聚类结果尤其需要警惕因命名而产生实体错觉。把某个簇称为``高价值用户''不会自动让它成为一个真实的自然类。在依赖聚类结果做决策之前，应至少检查：簇是否在不同随机子样本上稳定重现；簇的边界是否主要由少数几个特征的尺度差异驱动；新数据到来后能否被明确分配到一个簇；以及簇的差异是否足以支撑你计划采取的差异化行动。

K-means 提供的是一套清楚、可计算、可复现的分区准则。它发现的结构，精确地说，是``在当前特征表示和平方欧氏距离下，由 \(K\) 个中心较好解释的分组''。这后半句不附加到任何结论上，结论使用的范围就暗中被扩大到了准则未曾担保的领域。
